\documentclass[../../main.tex]{subfiles}
\graphicspath{{\subfix{../../image/}}}

\begin{document}

\section{上下解方法}\label{section:ch09-3-上下解方法}

\S9.2.2 中对 Schaefer 定理的应用依赖于椭圆方程解的正则性估计。现把注意力转向椭圆 PDE 的另一基本性质，即极值原理，并演示如何用由此导出的各种比较论证求解某些半线性问题。想法是利用解的序性质。更精确地，将证明：若能找到某边值问题的一个弱下解 $\underline{u}$ 与一个弱上解 $\overline{u}$，且 $\underline{u}\leqslant\overline{u}$，则事实上存在一个解 $u$ 满足
\begin{align*}
\underline{u}\leqslant u\leqslant\overline{u}.
\end{align*}

考察非线性 Poisson 方程的边值问题：
\begin{align}
\begin{cases}
-\Delta u=f(u)&\text{在 }U\text{ 中},\\
u=0&\text{在 }\partial U\text{ 上},
\end{cases}\label{align:ch09-3-eq01}
\end{align}
其中 $f:\mathbb{R}\to\mathbb{R}$ 光滑，且
\begin{align}
|f'|\leqslant C\qquad (z\in\mathbb{R})\label{align:ch09-3-eq02}
\end{align}
对某常数 $C$。

\begin{definition}[i]\label{definition:ch09-3-上下解} 称 $\overline{u}\in H^1(U)$ 是问题\eqref{align:ch09-3-eq01} 的弱上解，若
\begin{align}
\int_U D\overline{u}\cdot Dv\,\mathrm{d}x\geqslant\int_U f(\overline{u})v\,\mathrm{d}x\label{align:ch09-3-eq03}
\end{align}
对每个 $v\in H_0^1(U)$、$v\geqslant 0$ a.e. 成立。
(ii) 类似地，称 $\underline{u}\in H^1(U)$ 是弱下解，只要
\begin{align}
\int_U D\underline{u}\cdot Dv\,\mathrm{d}x\leqslant\int_U f(\underline{u})v\,\mathrm{d}x\label{align:ch09-3-eq04}
\end{align}
对每个 $v\in H_0^1(U)$、$v\geqslant 0$ a.e. 成立。
(iii) 称 $u\in H_0^1(U)$ 是\eqref{align:ch09-3-eq01} 的弱解，若
\begin{align*}
\int_U Du\cdot Dv\,\mathrm{d}x=\int_U f(u)v\,\mathrm{d}x
\end{align*}
对每个 $v\in H_0^1(U)$ 成立。
\end{definition}

\begin{remark}
若 $\overline{u},\underline{u}\in C^2(U)$，则由\eqref{align:ch09-3-eq03}、\eqref{align:ch09-3-eq04} 推得
\begin{align*}
-\Delta\overline{u}\geqslant f(\overline{u}),\qquad -\Delta\underline{u}\leqslant f(\underline{u})\quad\text{在 }U\text{ 中}.
\end{align*}
\end{remark}

\begin{theorem}[上下解之间解的存在性]\label{theorem:ch09-3-1-上下解之间存在}
设存在\eqref{align:ch09-3-eq01} 的弱上解 $\overline{u}$ 与弱下解 $\underline{u}$，满足
\begin{align}
\underline{u}\leqslant 0,\ \overline{u}\geqslant 0\text{ 在迹意义下于 }\partial U\text{ 上},\quad \underline{u}\leqslant\overline{u}\text{ a.e. 在 }U\text{ 中}.\label{align:ch09-3-eq05}
\end{align}
则存在\eqref{align:ch09-3-eq01} 的弱解 $u$，使
\begin{align*}
\underline{u}\leqslant u\leqslant\overline{u}\quad\text{a.e. 在 }U\text{ 中}.
\end{align*}
\begin{proof}
1. 取数 $\lambda>0$ 充分大使
\begin{align}
\text{映射 }z\mapsto f(z)+\lambda z\text{ 非减；}\label{align:ch09-3-eq06}
\end{align}
由假设\eqref{align:ch09-3-eq02}，这是可能的。现写 $u_0=\underline{u}$，然后给定 $u_k$（$k=0,1,2,\cdots$），归纳定义 $u_{k+1}\in H_0^1(U)$ 为线性边值问题
\begin{align}
\begin{cases}
-\Delta u_{k+1}+\lambda u_{k+1}=f(u_k)+\lambda u_k&\text{在 }U\text{ 中},\\
u_{k+1}=0&\text{在 }\partial U\text{ 上}
\end{cases}\label{align:ch09-3-eq07}
\end{align}
的唯一弱解。
2. 断言
\begin{align}
\underline{u}=u_0\leqslant u_1\leqslant\cdots\leqslant u_k\leqslant\cdots\quad\text{a.e. 在 }U\text{ 中}.\label{align:ch09-3-eq08}
\end{align}
为确认此点，先由\eqref{align:ch09-3-eq07}（$k=0$）注意
\begin{align}
\int_U Du_1\cdot Dv+\lambda u_1v\,\mathrm{d}x=\int_U(f(u_0)+\lambda u_0)v\,\mathrm{d}x\label{align:ch09-3-eq09}
\end{align}
对每个 $v\in H_0^1(U)$。由\eqref{align:ch09-3-eq04} 减去\eqref{align:ch09-3-eq09}，回顾 $u_0=\underline{u}$，并令 $v=(u_0-u_1)^+\in H_0^1(U)$、$v\geqslant 0$ a.e.。得
\begin{align}
\int_U D(u_0-u_1)\cdot D(u_0-u_1)^++\lambda(u_0-u_1)(u_0-u_1)^+\,\mathrm{d}x\leqslant 0.\label{align:ch09-3-eq10}
\end{align}
但
\begin{align*}
D(u_0-u_1)^+=\begin{cases}D(u_0-u_1)&\text{a.e. 在 }\{u_0\geqslant u_1\}\text{ 上},\\ 0&\text{a.e. 在 }\{u_0\leqslant u_1\}\text{ 上}.\end{cases}
\end{align*}
（见\cref{exercise:ch05-18}。）于是
\begin{align*}
\int_{\{u_0\geqslant u_1\}}|D(u_0-u_1)|^2+\lambda(u_0-u_1)^2\,\mathrm{d}x\leqslant 0,
\end{align*}
故 $u_0\leqslant u_1$ a.e. 在 $U$ 中。
现归纳假设
\begin{align}
u_{k-1}\leqslant u_k\quad\text{a.e. 在 }U\text{ 中}.\label{align:ch09-3-eq11}
\end{align}
由\eqref{align:ch09-3-eq07}，得
\begin{align}
\int_U Du_{k+1}\cdot Dv+\lambda u_{k+1}v\,\mathrm{d}x=\int_U(f(u_k)+\lambda u_k)v\,\mathrm{d}x,\label{align:ch09-3-eq12}
\end{align}
且
\begin{align*}
\int_U Du_k\cdot Dv+\lambda u_kv\,\mathrm{d}x=\int_U(f(u_{k-1})+\lambda u_{k-1})v\,\mathrm{d}x
\end{align*}
对每个 $v\in H_0^1(U)$。相减并令 $v=(u_k-u_{k+1})^+$。推得
\begin{align*}
\int_{\{u_k\geqslant u_{k+1}\}}|D(u_k-u_{k+1})|^2+\lambda(u_k-u_{k+1})^2\,\mathrm{d}x\\
=\int_U[(f(u_{k-1})+\lambda u_{k-1})-(f(u_k)+\lambda u_k)](u_k-u_{k+1})^+\,\mathrm{d}x\leqslant 0,
\end{align*}
最后一不等式由\eqref{align:ch09-3-eq11} 与\eqref{align:ch09-3-eq06} 成立。故 $u_k\leqslant u_{k+1}$ a.e. 在 $U$ 中，如所断言。
3. 接着证明
\begin{align}
u_k\leqslant\overline{u}\quad\text{a.e. 在 }U\text{ 中}\ (k=0,1,\cdots).\label{align:ch09-3-eq13}
\end{align}
由假设\eqref{align:ch09-3-eq05}，\eqref{align:ch09-3-eq13} 对 $k=0$ 成立。现归纳假设
\begin{align}
u_k\leqslant\overline{u}\quad\text{a.e. 在 }U\text{ 中}.\label{align:ch09-3-eq14}
\end{align}
然后由\eqref{align:ch09-3-eq12} 减去\eqref{align:ch09-3-eq03}，并取 $v=(u_{k+1}-\overline{u})^+$，得
\begin{align*}
\int_{\{u_{k+1}\geqslant\overline{u}\}}|D(u_{k+1}-\overline{u})|^2+\lambda(u_{k+1}-\overline{u})^2\,\mathrm{d}x\\
\leqslant\int_U[(f(u_k)+\lambda u_k)-(f(\overline{u})+\lambda\overline{u})](u_{k+1}-\overline{u})^+\,\mathrm{d}x\leqslant 0,
\end{align*}
由\eqref{align:ch09-3-eq14} 与\eqref{align:ch09-3-eq06}。故 $u_{k+1}\leqslant\overline{u}$ a.e. 在 $U$ 中。
4. 由\eqref{align:ch09-3-eq08} 与\eqref{align:ch09-3-eq13}，有
\begin{align}
\underline{u}\leqslant\cdots\leqslant u_k\leqslant u_{k+1}\leqslant\cdots\leqslant\overline{u}\quad\text{a.e. 在 }U\text{ 中}.\label{align:ch09-3-eq15}
\end{align}
因此
\begin{align*}
u(x)\triangleq\lim_{k\to\infty}u_k(x)
\end{align*}
对 a.e. $x$ 存在。且由控制收敛定理与\eqref{align:ch09-3-eq15}，
\begin{align}
u_k\to u\quad\text{于 }L^2(U).\label{align:ch09-3-eq16}
\end{align}
最后，因 $\|f(u_k)\|_{L^2(U)}\leqslant C(\|u_k\|_{L^2(U)}+1)$，由\eqref{align:ch09-3-eq07} 推得 $\sup_k\|u_k\|_{H_0^1(U)}<\infty$。故存在子列 $\{u_{k_j}\}_{j=1}^\infty$ 在 $H_0^1(U)$ 中弱收敛到 $u\in H_0^1(U)$。
5. 最后验证 $u$ 是问题\eqref{align:ch09-3-eq01} 的弱解。为此固定 $v\in H_0^1(U)$。由\eqref{align:ch09-3-eq07} 得
\begin{align*}
\int_U Du_{k+1}\cdot Dv+\lambda u_{k+1}v\,\mathrm{d}x=\int_U(f(u_k)+\lambda u_k)v\,\mathrm{d}x.
\end{align*}
令 $k\to\infty$：
\begin{align*}
\int_U Du\cdot Dv+\lambda uv\,\mathrm{d}x=\int_U(f(u)+\lambda u)v\,\mathrm{d}x.
\end{align*}
消去含 $\lambda$ 的项，最终确认
\begin{align*}
\int_U Du\cdot Dv\,\mathrm{d}x=\int_U f(u)v\,\mathrm{d}x,
\end{align*}
如所要求。
\end{proof}
\end{theorem}

此例证了用分部积分而非极值原理来建立上下解之间比较。

\end{document}
